% !TEX root = ./index.tex

\maketitle
\newpage

\section*{Answers}

\bigskip

\begin{enumerate}

        % 1.
    \item
        \begin{enumerate}

                % 1. a.
            \item Let $U$ be a set and let $C \subseteq U$ be fixed.
                Define
                \[ f_C: \mathcal{P}(U) \to \mathcal{P}(U) \times \mathcal{P}(U), \]
                \[ f_C(A) = \left(A \cap C, A \cup C\right). \]

                \medskip

                To prove the following proposition, we shall use the following lemma:
                \begin{lemma} \label{lem:1a-1}
                    For any set $K$, the set
                    \[ S = (K \cap L) \cup \left((K \cup L) \, \backslash \, L\right) \]
                    is equivalent to $K$.
                \end{lemma}

                \medskip

                \begin{proof}
                    For any element $x$ in $K$, either $x \in L$ or $x \notin L$.

                    If $x \in L$, then it is also a member of $K \cap L$ since $x \in K$ and $x \in L$.
                    If $x \notin L$, then it is a member of $(K \cup L) \, \backslash \, L$ since $x \in K$ and $x \notin L$.
                    Either way, $x$ is a member of $S$ and hence $K \subseteq S$.

                    \smallskip

                    Note that the sets $(K \cup L) \, \backslash \, L$ and $K \cap L$ are both subsets of $K$.
                    Since $S$ is the union of these two subsets, it follows that $S \subseteq K$.

                    We have $K \subseteq S$ and $S \subseteq K$, hence $S = K$.
                \end{proof}

                \medskip

                \begin{proposition} \label{prop:1a}
                    The function $f_C$ is injective.
                \end{proposition}

                \medskip

                \begin{proof}
                    Take two arbitrary subsets $X$ and $Y$ of $U$.
                    If $f_C(X) = f_C(Y)$, then we have
                    \begin{equation} \label{eq:1a-1}
                        X \cap C = Y \cap C, \qquad X \cup C = Y \cup C.
                    \end{equation}

                    Substituting $K = X$, $L = C$ into \cref{lem:1a-1},we obtain the equality
                    \begin{align*}
                        X &= (X \cap C) \cup ((X \cup C) \, \backslash \, C). \\
                        \intertext{Similarly,}
                        Y &= (Y \cap C) \cup ((Y \cup C) \, \backslash \, C).
                    \end{align*}
                    From \cref{eq:1a-1}, it follows that $X = Y$, therefore
                    \[ f_C(X) = f_C(Y) \quad \implies \quad X = Y \]
                    and $f_C$ is injective.
                \end{proof}

                \bigskip

                % 1. b.
            \item
                \begin{proposition} \label{prop:1b}
                    The range of $f_C$ is
                    \[ \mathrm{Range}(f_C) = \left\{(X, Y) \; \big| \; X \subseteq C \subseteq Y \subseteq U\right\}. \]
                \end{proposition}

                \medskip

                \begin{proof}
                    The ordered pair $(X, Y)$ is in the range of $f_C$ if and only if
                    \[ X = A \cap C, \qquad Y = A \cup C \]
                    for some $A \subseteq U$.

                    Suppose that this is true.
                    Then by \cref{lem:1a-1},
                    \[ A = X \cup (Y \, \backslash \, C). \]
                    Note that $A \subseteq U$, so $X \subseteq U$ and $Y \subseteq U$.

                    We can substitute this expression into the expressions for $X$ and $Y$ to get
                    \begin{align*}
                        X &= \left[X \cup (Y \, \backslash \, C)\right] \cap C, \\
                        &= (X \cap C) \cup \left[(Y \, \backslash \, C) \cap C\right], \\
                        &= (X \cap C) \cup \varnothing, \\
                        &= X \cap C, \\
                        \intertext{and}
                        Y &= \left[X \cup (Y \, \backslash \, C)\right] \cup C, \\
                        &= X \cup \left[(Y \, \backslash \, C) \cup C\right], \\
                        &= X \cup (Y \cup C).
                    \end{align*}

                    The equation $X = X \cap C$ implies that $X \subseteq C$.
                    With this in mind, the equation $Y = X \cup (Y \cup C)$ simplifies to $Y = Y \cup C$, which implies that $Y \supseteq C$.

                    Hence if $(X, Y)$ is in the range of $f_C$, then $X \subseteq C \subseteq Y \subseteq U$.
                    We write
                    \[ \mathrm{Range}(f_C) \subseteq \left\{(X, Y) \; \big| \; X \subseteq C \subseteq Y \subseteq U \right\}. \]

                    \smallskip

                    To prove equality, fix $S$ and $T$ such that $S \subseteq C \subseteq T \subseteq U$, and set
                    \[ A = S \cup (T \, \backslash \, C). \]
                    By the logic described above, computing the intersection and union of $A$ with $C$ gives
                    \begin{align*}
                        A \cap C &= S, \\
                        A \cup C &= T.
                    \end{align*}
                    Thus, there exists $A \in U$ such that $f_C(A) = (S, T)$. We write
                    \[ \mathrm{Range}(f_C) \supseteq \left\{(X, Y) \; \big| \; X \subseteq C \subseteq Y \subseteq U \right\}, \]
                    hence by the definition of equality,
                    \[ \mathrm{Range}(f_C) = \left\{(X, Y) \; \big| \; X \subseteq C \subseteq Y \subseteq U \right\}. \qedhere \]
                \end{proof}

                %TODO 1. c.
            \item Define $g, h: \mathcal{P}(U) \to \mathcal{P}(U)$ by
                \[ g(A) = A \cap C, \qquad h(A) = A \cup C. \]

                \begin{proposition} \label{prop:1c}
                    $g$ is injective if and only if $C = U$, and $h$ is injective if and only if $C = \varnothing$.
                \end{proposition}

                \medskip

                \begin{proof}
                    It is relatively simple to prove that if $C = U$, then $g$ is injective.
                    Since $A \subseteq U$, we have
                    \[ g(A) = A \cap U = A, \]
                    which is the injective identity function.

                    Likewise, when $C = \varnothing$,
                    \[ h(A) = A \cup \varnothing = A, \]
                    which is again the identity function.

                    \medskip

                    We now want to prove the reverse implication, or rather its contrapositive, beginning with $g$:
                    \begin{center}
                        If $C \neq U$, then $g$ is not injective.
                    \end{center}

                    Assume that $C \neq U$, and consider the values of $g(U)$ and $g(C)$.
                    Since $C \subset U$, we have
                    \[ g(U) = U \cap C = C, \]
                    and
                    \[ g(C) = C \cap C = C. \]
                    Hence $g(U) = g(C)$, which contradicts the requirement for injectivity.

                    \smallskip

                    A similar method is used to prove the equivalent statement for $h$.

                    Assume that $C \neq \varnothing$, and consider the values of $h(\varnothing)$ and $h(C)$.
                    Since $C \supset \varnothing$, we have
                    \[ h(\varnothing) = \varnothing \cup C = C, \]
                    and
                    \[ h(C) = C \cup C = C. \]
                    Hence $h(\varnothing) = h(C)$, which again contradicts the requirement for injectivity.
                \end{proof}

                %TODO 1. d.
            \item

        \end{enumerate}

        \setcounter{enumi}{3}

        %TODO 4.
    \item
        \begin{enumerate}

                %TODO 1. a.
            \item

                %TODO 1. b.
            \item

                %TODO 1. c.
            \item

                %TODO 1. d.
            \item

        \end{enumerate}

\end{enumerate}
